\section{Photovoltaic System Modeling}
\label{sec:pv-model}

\subsection{Two-Diode Model}
The reference module is a Canadian Solar CS6P-250P, a 60-cell polycrystalline
module rated at $P_{max}=250\,$W ($V_{oc}=37.2\,$V, $I_{sc}=8.87\,$A,
$V_{mp}=30.1\,$V, $I_{mp}=8.30\,$A at standard test conditions, STC:
1000\,W/m$^2$, 25\textdegree C), chosen as a real, commercially available
panel rather than an idealized one so simulated results can be validated
quantitatively against the manufacturer's own datasheet.

Each series-connected cell string is modeled by the two-diode equivalent
circuit
\begin{equation}
I = I_{ph} - I_{s1}\!\left[e^{\frac{V+IR_s}{a_1 V_t}} - 1\right]
          - I_{s2}\!\left[e^{\frac{V+IR_s}{a_2 V_t}} - 1\right]
          - \frac{V+IR_s}{R_{sh}},
\label{eq:two-diode}
\end{equation}
where $V_t = a\,N_s\,kT/q$ is the string-level thermal voltage, $N_s$ the
number of series cells in the string, $k$ Boltzmann's constant, $T$ the
cell temperature in Kelvin, and $q$ the electron charge. Unlike the
single-diode model, the second diode term (with ideality $a_2$ and
saturation current $I_{s2}$) explicitly represents recombination current in
the depletion region, which becomes significant at low forward bias -- the
regime relevant to correctly resolving secondary local maxima under partial
shading. We follow \cite{ishaque2011simple} for the two-diode formulation;
\cite{villalva2009comprehensive}'s widely-cited three-point extraction
method is a single-diode technique and is cited here only for general
parameter-extraction background, not as the source of
\eqref{eq:two-diode}.

\subsection{Bypass Diodes and Partial Shading}
Following standard industry practice, one bypass diode protects every 20
cells, giving $N_s/20=3$ bypass groups for the 60-cell reference module.
Each group is modeled as an independent two-diode string of
\eqref{eq:two-diode}; when a forced string current would exceed a group's
own short-circuit current at its (possibly shaded) irradiance, that group's
bypass diode is modeled as an ideal diode with $V_f=0.7\,$V, clamping the
group's contributed voltage to $-V_f$ rather than driving it into deep
reverse bias. This bypass mechanism is what allows multiple local maxima to
appear in a shaded string's $P$-$V$ curve at all -- without it, a
series-connected string's power curve remains unimodal regardless of
shading, which would make the partial-shading scenarios in
Section~\ref{sec:scenarios} unable to test any algorithm's global-search
capability.

\subsection{Parameter Extraction}
\label{sec:parameter-extraction}
The four free parameters $(I_{ph}, R_s, R_{sh}, a_2)$ are fit to the
datasheet's four STC key points ($V_{oc}$, $I_{sc}$, $V_{mp}$, $I_{mp}$) via
bounded nonlinear least squares (Trust Region Reflective), minimizing the
residual of three conditions: the short-circuit current match, the
current-at-$V_{mp}$ match, and $dP/dV=0$ at the reported MPP. The first
diode's ideality factor $a_1$ is fixed at $1.0$ (the midpoint of its
physically expected range) rather than fit, and $I_{s1}=I_{s2}$ is solved
analytically from the open-circuit-voltage boundary condition at every
optimizer iteration rather than treated as a free parameter. This leaves
the extraction under-determined by one degree of freedom relative to the
full iterative procedure in \cite{ishaque2011simple}; we flag this
explicitly as a simplification rather than presenting it as the complete
published procedure.

For the CS6P-250P, this extraction yields $I_{ph}=8.880\,$A,
$I_{s1}=I_{s2}=2.891\times10^{-10}\,$A, $R_s=0.308\,\Omega$,
$R_{sh}=269.8\,\Omega$, $a_1=1.000$, $a_2=1.406$.

\subsection{Model Validation}
\label{sec:model-validation}
Table~\ref{tab:validation} compares the simulated I-V curve at the four STC
key points against the datasheet. The maximum relative error across all
four points is $0.068\%$ (at $P_{max}$), well within the $2\%$ acceptance
criterion.

\begin{table}[htbp]
\centering
\caption{Simulated vs.\ Datasheet STC Key Points (CS6P-250P)}
\label{tab:validation}
\begin{tabular}{@{}lccc@{}}
\toprule
Quantity & Simulated & Datasheet & Rel.\ Error \\
\midrule
$I_{sc}$ (A)  & 8.870  & 8.870  & 0.000\% \\
$V_{oc}$ (V)  & 37.200 & 37.200 & 0.000\% \\
$V_{mp}$ (V)  & 30.100 & 30.100 & 0.000\% \\
$I_{mp}$ (A)  & 8.300  & 8.300  & 0.000\% \\
$P_{max}$ (W) & 249.83 & 250.00 & 0.068\% \\
\bottomrule
\end{tabular}
\end{table}

\subsection{Temperature Dependence}
\label{sec:temperature-dependence}
Photocurrent and both diodes' saturation currents are temperature-dependent,
following \cite{desoto2006improvement}:
\begin{align}
I_{ph}(G,T) &= I_{ph,ref}\,\frac{G}{G_{ref}}\bigl(1 + K_i(T-T_{ref})\bigr), \label{eq:iph-temp}\\
I_{s}(T) &= I_{s,ref}\left(\frac{T}{T_{ref}}\right)^{3/a}
  \exp\!\left[\frac{E_g q}{ak}\left(\frac{1}{T_{ref}}-\frac{1}{T}\right)\right], \label{eq:is-temp}
\end{align}
applied independently to $I_{s1}$ (using $a_1$) and $I_{s2}$ (using $a_2$),
with $E_g=1.121\,$eV the silicon bandgap energy and $K_i=+0.065\%/^\circ$C
the datasheet's short-circuit-current temperature coefficient. Both
correction factors in \eqref{eq:iph-temp}-\eqref{eq:is-temp} are identically
1 at $T=T_{ref}$, so STC behavior is unaffected by this extension --
Table~\ref{tab:validation} above is unchanged whether or not temperature
dependence is included.

This addresses a real error found during development: an earlier version of
this model let temperature affect only the thermal-voltage term $V_t(T)$,
leaving $I_{ph}$ and $I_s$ temperature-independent. Holding $I_s$ fixed while
$V_t$ rises with temperature makes $V_{oc}\approx V_t\ln(I_{ph}/I_s)$
\emph{increase} with temperature -- backwards from every real PV panel,
where $I_s$'s much stronger (exponential) growth with temperature is the
physically dominant effect and pulls $V_{oc}$ \emph{down}. Concretely, the
uncorrected model predicted $V_{oc}(50^\circ\text{C})=40.32\,$V against this
panel's datasheet-implied $34.04\,$V (using its linear
$K_v=-0.34\%/^\circ$C coefficient) -- the wrong sign, not merely an
uncalibrated magnitude. \eqref{eq:iph-temp}-\eqref{eq:is-temp} fix this:
the corrected model gives $V_{oc}(50^\circ\text{C})=34.31\,$V, within
$0.8\%$ of the datasheet-implied value, and $I_{sc}(50^\circ\text{C})$
matches to within $0.01\%$ (the \emph{temperature step} scenario in
Section~\ref{sec:results} is the only one of this paper's eleven scenarios
that varies temperature -- every other scenario runs at a constant
25\textdegree C and is completely unaffected by this fix).

Table~\ref{tab:noct-validation} validates the corrected model against the
manufacturer's NOCT electrical table ($G=800\,$W/m$^2$, $T_c=45\pm2^\circ$C)
-- data not available until the full datasheet PDF was sourced during this
project. RMSE across the five NOCT key points is $1.35\%$, within the $2\%$
acceptance criterion (the individual $P_{max}$ point alone is $2.44\%$,
narrowly over 2\% on its own, though the aggregate RMSE the criterion
specifies is not). Low-irradiance validation is necessarily coarser: the
datasheet gives only a single retention claim ("+95.5\% module efficiency
from 1000 to 200\,W/m$^2$"), not a full I-V table at 200\,W/m$^2$, so this
is a single ratio check rather than a point-by-point match like NOCT. The
model measures $93.5\%$ retention -- close to, but about 2 percentage
points below, the datasheet's claim, plausibly because this simplified
model holds $R_s$ and $R_{sh}$ irradiance-independent rather than scaling
$R_{sh}$ with irradiance as a full five-parameter model would.

\begin{table}[htbp]
\centering
\caption{NOCT Validation ($G=800$\,W/m$^2$, $T_c=45^\circ$C)}
\label{tab:noct-validation}
\begin{tabular}{@{}lccc@{}}
\toprule
Quantity & Simulated & Datasheet & Rel.\ Error \\
\midrule
$I_{sc}$ (A)  & 7.188  & 7.190  & 0.024\% \\
$V_{oc}$ (V)  & 34.514 & 34.200 & 0.917\% \\
$V_{mp}$ (V)  & 27.814 & 27.500 & 1.143\% \\
$I_{mp}$ (A)  & 6.666  & 6.600  & 0.999\% \\
$P_{max}$ (W) & 185.41 & 181.00 & 2.435\% \\
\bottomrule
\multicolumn{4}{l}{\footnotesize RMSE across all five points: 1.35\%.}
\end{tabular}
\end{table}
